\section{Certified arithmetic}\label{sec:rounding}

This section supplies the one-sided arithmetic estimates used in
Sections~\ref{sec:seed}, \ref{sec:completion}, \ref{sec:phase-cover},
\ref{sec:refinement}, and~\ref{sec:exclusion}.
The source files, parameters, build, and completed execution are
documented in Section~\ref{sec:execution}.  The estimates concern those entry points
and parameters.  Auxiliary routines not called by those entry points are
not additional exclusion tests.

We distinguish an \emph{exact stored number}, meaning the real number
represented by a binary64 datum, from the ideal mathematical quantity it
approximates.  Proposed solve directions, proposed inverse matrices, and
proposed changes of basis are usually used in the first sense.  Their
accuracy is not assumed: finite bounds and residual inequalities are
checked before a proposal can justify an exclusion or a covering step.
For complex arrays, an entrywise error is an error in complex modulus.
Matrix norms without a subscript are Euclidean operator norms.

\subsection{Arithmetic premises and elementary estimates}
\label{round:premises}

The primitive-operation model is the binary64 specification of
IEEE~754~\cite[Sections~3.4, 3.6, 4.3, 5.4.1 and~7.5]{IEEE7542019}.
The accumulation estimates use the usual product-of-rounding-factors
argument~\cite[Section~2.12 and Lemma~3.1]{Higham2002Accuracy}, with the
explicit additive underflow allowance retained below.

The arithmetic premises are separate, correctly rounded binary64
addition, subtraction, multiplication, division, and square root in
round-to-nearest mode, with binary64 evaluation of intermediate results.
The build uses \texttt{-ffp-contract=off}, and does not use fast-math or
reassociation that changes these operations.  Subnormal numbers are not
flushed to zero, and denormal inputs are not treated as zero.  Floating
exceptions are masked, so an unsuccessful heuristic calculation can
produce an infinity or NaN and be rejected by its subsequent checks.
The relevant integer types, numeric locale, and file representation are
specified in Section~\ref{sec:execution}.

The function \texttt{environment} sets and checks nearest rounding and
checks the binary64 properties of \texttt{double}.  The A/B/C entry points
also inspect MXCSR, rejecting a non-nearest mode, DAZ, FTZ, or unmasked
exceptions.  The same no-DAZ/no-FTZ environment is required for
\texttt{covout}, although its entry point does not itself inspect those
bits.  The default-nearest premise also applies to constants and the
global basis constructed before \texttt{main}.  These are explicit
arithmetic premises, not a claim that a source-level static assertion
formally verifies every compiler transformation.  The recorded build
checks and their limits are described in Section~\ref{sec:execution}.

The functions \texttt{nearbyint} and \texttt{floor} choose integers at
finite bounded arguments.  The particular integer chosen at a boundary
is permissible wherever the proof subsequently uses that integer lift.
No accuracy premise is imposed on \texttt{atan2}: Candidate checks the
unit exponential of the phase it proposes.  Decimal literals are either
used as exact stored parameters, or allowed a relative translation error
of at most $5\cdot10^{-16}$ when compared with a displayed decimal.

Put
\begin{equation}\label{round:unit}
 u=2\cdot10^{-16},\qquad \tau=10^{-300},\qquad
 \gamma_m=2.1\cdot10^{-16}(m+1).
\end{equation}
For a correctly rounded basic operation with finite exact result $z$ at
the scales below,
\begin{equation}\label{round:basic}
 |\operatorname{fl}(z)-z|\le u|z|+\tau.
\end{equation}
The true binary64 relative unit roundoff is smaller than $u$, and the
additive allowance is much larger than half the least positive
subnormal.  Overflows are excluded for the certificate evaluations by
the scale bounds proved below; they need not be excluded in an
untrusted solve.

\begin{lemma}[Accumulation and norm bounds]\label{round:accumulation}
For $m\le2000$, a dot product on exact stored inputs, formed as separate
products and successive additions, satisfies
\begin{equation}\label{round:dot}
 \left|\operatorname{fl}\!\left(\sum_{i=1}^m a_i b_i\right)
                  -\sum_{i=1}^m a_i b_i\right|
 \le \gamma_m\sum_{i=1}^m|a_i b_i|+4(m+1)\tau.
\end{equation}
An initial addend is included in the absolute sum and increases the
small operation-count constant by one.  A product of three factors
instead of two increases that constant by one more.
For the square-sum/square-root norms used here,
\begin{equation}\label{round:norm}
 \left|\operatorname{fl}(\|a\|)-\|a\|\right|
 \le 2.2\cdot10^{-16}(m+2)\|a\|+10^{-140}.
\end{equation}
If the computed components approximate other components, their
Euclidean difference is added to this bound.
\end{lemma}

\begin{proof}
For each product term, the product rounding and the additions through
which it passes contribute at most $m+1$ multiplicative factors of the
form $1+\delta$, $|\delta|\le u$.  In the stated range,
$(1+u)^{m+1}-1<\gamma_m$.  Propagating the additive errors in
\eqref{round:basic} costs less than $4(m+1)\tau$; all amplification
factors here are less than $1.001$.  This proves \eqref{round:dot}.
The same argument applies to positive sums and to three-factor products.
For a sum of squares it gives relative error at most $\gamma_m$ and an
additive error at most $4(m+1)\tau$.  Taking its square root, using
$|\sqrt{1+t}-1|\le |t|/(1+\sqrt{1-|t|})$ for $|t|<1$, and then including
the final square-root rounding gives \eqref{round:norm}.  The additive
effect is bounded by $\sqrt{4(m+1)\tau}$ plus the square-root rounding;
this is smaller than $10^{-140}$.  Finally
$|\|a\|-\|b\||\le\|a-b\|$.  Thus no division by a small norm is needed
to handle perturbed components.
\end{proof}

For exact stored complex numbers, the explicit source operations obey
\begin{align}
 |\operatorname{fl}(z+w)-(z+w)|
   &\le u(|z|+|w|)+2\tau,\label{round:complex-add}\\
 |\operatorname{fl}(zw)-zw|
   &\le 10^{-15}|z||w|+8\tau.\label{round:complex-multiply}
\end{align}
For the latter, each real component uses two real products and an
addition or subtraction.  Its error is bounded by
$(2u+u^2)$ times the sum of the two absolute products, plus the additive
allowance.  Cauchy's inequality bounds each such sum by $|z||w|$, and
$\sqrt2(2u+u^2)<10^{-15}$.  Real scaling and division are componentwise
instances of \eqref{round:basic}.  Squared moduli also satisfy the useful
perturbation bound
\begin{equation}\label{round:squared-perturbation}
 \big||z+e|^2-|z|^2\big|\le2|z||e|+|e|^2,
\end{equation}
before their own multiplication/addition errors.

The additive underflow terms can be retained throughout these
inequalities.  The norm allowance $10^{-140}$ is deliberately loose.
All subsequent certified divisions have the explicit lower bounds
given below; in particular the residual test divides only by a gap
of order $10^{-6}$ or larger.  Its norm factors and scale are at most
the displayed $10^8$--$10^{11}$ scales, so its propagated additive
terms remain below $10^{-120}$.  The second square root in the spectral
estimate is handled separately by an absolute term of order
$10^{-69}$.  That term is covered there before the result is reused.
Other normalizations have denominators bounded below by $.148$, $.45$,
or larger.  At each iterative upper-bound update the new local error is
absorbed by its own positive margin.  There is therefore no unchecked
chain of ill-conditioned divisions propagating a subnormal error: all
such additive contributions are below $10^{-50}$ when they are
absorbed in the strict budgets below.

\subsection{Pi and unit exponentials}
\label{round:exponentials}

The source literal \texttt{0x1.921fb54442d18p+1} represents the
binary rational
\[
 \widehat\pi=\frac{884279719003555}{281474976710656}.
\]
Here is an exact certificate for its error.  Write
\[
 A(x,n)=\sum_{j=0}^{n-1}\frac{(-1)^j x^{2j+1}}{2j+1},\qquad
 s=16A(1/5,12)-4A(1/239,4).
\]
If $t=\arctan(1/5)$, then $\tan(2t)=5/12$ and
$\tan(4t)=120/119$; subtracting $\arctan(1/239)$ gives tangent $1$
in $(0,\pi/2)$.  This proves Machin's identity
$\pi=16\arctan(1/5)-4\arctan(1/239)$.
The alternating-series remainder gives
\[
 |s-\pi|\le \frac{16}{25\,5^{25}}+\frac{4}{9\,239^9}.
\]
Exact rational arithmetic gives
\begin{equation}\label{round:pi-rational}
 \widehat\pi-s=
 \frac{-1204334632739641172855772829555432415587}
 {10003273510806667494906565283020800000000000000000000000}.
\end{equation}
Adding the displayed rational remainder bounds and comparing integers
proves
\begin{equation}\label{round:pi-error}
 |\widehat\pi-\pi|<1.225418\cdot10^{-16}<2\cdot10^{-16}.
\end{equation}

The range-reduction argument uses Sterbenz's exact-subtraction lemma
\cite{Sterbenz1974FloatingPoint}; a statement including gradual
underflow is given in~\cite[Theorem~2.5 and the following discussion]{Higham2002Accuracy}.
The polynomial and argument-error estimates are proved here for the
particular degree-$18/19$ implementation.

\begin{lemma}[Unit-exponential error]\label{round:expair}
At a finite stored $x$ with $|x|\le16$, \texttt{expair} differs from
$e(x)=\exp(2\pi i x)$ in complex modulus by less than $2\cdot10^{-14}$.
For all its uses in the certificate, including argument formation,
the following common bound is valid:
\begin{equation}\label{round:exp-budget}
 \varepsilon_{\exp}=6\cdot10^{-14}.
\end{equation}
\end{lemma}

\begin{proof}
The multiplication $4x$ is exact.  Let $k$ be its nearest integer.
Then $|x-k/4|\le1/8$.  If $k\ne0$, $x$ and $k/4$ have the same sign
and their absolute values are within a factor two, so the subtraction
is exact by Sterbenz's lemma.  If $k=0$ it is immediate.  Multiplication
by the exact double $2\widehat\pi$ forms a stored angle $y$ with
\[
 |y-2\pi(x-k/4)|<3\cdot10^{-16},\qquad |y|<.786.
\]
The error follows from \eqref{round:pi-error}, the remainder bound $1/8$,
and one multiplication rounding.

The source evaluates cosine through degree 18 and sine through degree
19.  In the $j$th nonconstant term, the error in $-y^2$ is reused $j$
times, and there are $j$ further multiplications and $j$ divisions.
The small integer denominators are exact.  Its relative error is
therefore at most $(1+u)^{3j}-1<5.5\cdot10^{-15}$ for $j\le9$.
For additive underflow errors, compare with a recurrence using the
same multiplicative rounding factors but zero additive errors.
Its terms have magnitude below two, and each recurrence contracts
the previous additive error by a factor less than $1/2$:
the numerator factor has magnitude below $.786^2$ and the smallest
denominator is two.  Errors from that factor, the product, and the
division give the conservative recurrence
$E_j\le E_{j-1}/2+4\tau$, $E_0=0$.
Thus $E_j<8\tau$; after both coordinates' nine additions their
total contribution is below $10^{-295}$, including when $y^2$
or a later term underflows.
The absolute even-term sum at $.786$ is below $1.4$, its nonconstant
part below $.4$; the odd-term sum is below $1$, its part after $y$
below $.2$.  For the nonconstant even terms the successive ratio
after the first is at most $.786^2/12$, and for the odd terms after y
it is at most $.786^2/20$.  Their geometric majorants give the claimed
absolute sums without evaluating a transcendental function.
The first omitted cosine and sine terms
are below $3.330\cdot10^{-21}$ and $1.247\cdot10^{-22}$, respectively;
in particular both tails are below $4\cdot10^{-20}$.

The term errors contribute less than $2.2\cdot10^{-15}$ in cosine and
$1.1\cdot10^{-15}$ in sine.  Put
$\alpha=(1+u)^9-1<1.81\cdot10^{-15}$ and
$\rho=(1+u)^{27}-1<5.5\cdot10^{-15}$.
Including the nine additions and the tails gives coordinate errors
less than
$.4\rho+1.4(1+\rho)\alpha+4\cdot10^{-20}<4.74\cdot10^{-15}$
and
$.2\rho+(1+\rho)\alpha+4\cdot10^{-20}<2.92\cdot10^{-15}$.
Their Euclidean combination, the $3\cdot10^{-16}$ angle error,
and the additive underflow term give complex error below
$6\cdot10^{-15}$, hence below the asserted
$2\cdot10^{-14}$.  The quadrant permutations and sign changes are exact.

For a rational bin midpoint, division by the grid size, or multiplication
of $b+1/2$ by the stored reciprocal of that size, has error below
$5\cdot10^{-16}$ cycles.  The latter product has ideal magnitude at
most one.  Subtraction of a phase of magnitude at most four, or a
stored ball phase of magnitude at most two, brings the total argument
error to less than $2\cdot10^{-15}$ cycles.  The Lipschitz constant of
$e$ is $2\pi$, so this additional error fits inside
\eqref{round:exp-budget}.  Dyadic unitary centers are exact; Stage A
treats the stored seed phases as exact after enlarging the seed ball;
\texttt{buildH} treats its stored $p$ as exact and compares row-edge
exponentials with products of ideal unit entries.
\end{proof}

All relied-on arguments satisfy the range condition.  Seed and dyadic
unitary phases are bounded near $[0,1]$; \texttt{gen} checks $|p_j|\le4$;
Candidate checks $|P_j|\le.51$ before its exponential; C validates
$|p_j|\le4$ and checks ball phases of magnitude at most two.
The largest C difference has magnitude at most six.  The constants
$\sqrt6$ and its reciprocal are formed by the indicated square root
and division; $2\cdot10^{-15}$ is an ample absolute error bound for each.

\subsection{Seed filtering and seed-ball reuse}
\label{round:seed}

Here $n=112$ and the stored $w$ approximates $1/n$.
Midpoint-gap calculations contain only small integer arithmetic and
a few products/subtractions; \eqref{round:basic} bounds their total
error, including the comparison threshold, by $10^{-14}$.
Their slack is $10^{-8}$.  The orientation filter is integral.

For a list of five varying phases, the computed residual $c$ in
\texttt{sumok} has component error below $2\cdot10^{-13}$.
Indeed its five exponential errors, after division by $2\pi$, total
less than $4.8\cdot10^{-14}$ in complex modulus, and the divisions,
additions, and constant error cost less than $3\cdot10^{-15}$.
The sharper componentwise bound $1.1\cdot10^{-13}$ will also be useful.
Each test direction is regarded as an exact stored pair; its norm is
less than $1.1$.  The five column perturbations, the residual
perturbation, and ordinary arithmetic in \texttt{chk} cost together
less than $10^{-11}$.  Its remainder coefficient approximates
$5\pi/(4n^2)$ within $10^{-15}$.  The actual $10^{-8}$ comparison
addition therefore preserves every feasible list.

For \texttt{extend}, the stored direction has norm at most $1.01$.
Let $v$, $a$, and $r$ denote its stored directional residual in grid
units, shared coefficient, and support before inflation.  Their errors
satisfy
\begin{equation}\label{round:clip-errors}
 |\Delta v|<3.3\cdot10^{-11},\qquad
 |\Delta a|<10^{-13},\qquad |\Delta r|<10^{-12}.
\end{equation}
For the first bound, multiply the sharpened residual bound by $112$
and by a direction of norm at most $1.01$, and include a two-term dot
and one multiplication.  The other two follow from
\eqref{round:exp-budget}, four absolute support terms, and a
two-component norm.  If $|d_s|\le1/2$ is feasible, these estimates give
\[
 |v+a d_s|\le r+3.3\cdot10^{-11}
                 +\tfrac12\cdot10^{-13}+10^{-12}+\text{rounding}.
\]
The source adds $10^{-7}$ to $r$ \emph{before} dividing by $|a|$.
It skips $|a|<10^{-6}$; at an accepted divisor
$|a|>.999\cdot10^{-6}$, allowing literal translation.
The midpoint quotient has magnitude below $1.2\cdot10^8$ and the
radius quotient below $3\cdot10^6$: the residual has magnitude below
112, and the support before division is below three.  Their divisions
and endpoint additions can move an endpoint by less than $10^{-7}$.
The extra $10^{-5}$ on each endpoint is therefore sufficient.
Intersecting these outward intervals retains the true shared value.

In \texttt{supdual}, all proposed entries are checked by
$|v_a|\le10^4$, with a negated comparison that also rejects NaN.
They are then exact stored test coefficients.  The four right-hand
side errors against $-nc$ are below $3\cdot10^{-11}$ each.  The signed
term consequently has input error at most
\[
 2\sum_{a=1}^4 |v_a|\,3\cdot10^{-11}
 \le2.4\cdot10^{-6}<2.5\cdot10^{-6}.
\]
Its absolute term sum is below $10^7$, so its arithmetic error is below
$10^{-7}$.  Each of the nine row-dot residuals has error below
$4\cdot10^{-9}$: its four input contributions are at most
$4\cdot10^4\varepsilon_{\exp}$, and its short dot costs less than
$10^{-10}$.  Each pair norm has error below $10^{-10}$ by
\eqref{round:norm}.  The remainder coefficient uses $1.251$ instead
of $1.25$, which leaves positive slack after its own rounding.
The final absolute scale is below $2\cdot10^7$.  Summing these
contributions shows that the ideal support expression is at most the
computed proposal plus $10^{-4}$.  The actual addition is $.01$.
If no usable proposal improves the initial $10^{10}$ bound, that
initial bound already exceeds the trivial support.

Integer differences lie in $[-111,111]$, including those stored in
\texttt{int8\_t}; all their squared/absolute sums are exact and small.
The comparison radius is
\[
 \texttt{RN2}=\operatorname{fl}((\operatorname{fl}(.032\cdot112))^2)
              -.0001
\]
with the source's indicated operation ordering.  The total possible
error before the subtraction is below $10^{-12}$, so this is smaller
than the ideal $(.032\cdot112)^2$.  This reasoning requires the
\texttt{pow(x,2)} calls in \texttt{covout} to be lowered to binary64
multiplication, as checked for the fresh build documented in
Section~\ref{sec:execution}.  It is not a claim about an arbitrary
library implementation of \texttt{pow}.  The other square-power value
$T$, bucket reach, and integer cuts only select reuse attempts; missing
an attempt does not discard the current bin pair.  A newly inserted
midpoint needs radius only $1.5/112<.032$.  Finally, the nine
division-rounded seed phases differ from their rational midpoints in
norm by less than $10^{-14}$, which is absorbed when the A seed radius
is increased to $.0320001$.

\subsection{Frames, comparison blocks, and A exclusion tests}
\label{round:frames}

All ideal quantities in this subsection use exact unit exponentials
at the specified centers, and the same coordinate indices selected by
the code.  We never need the selected coordinate to maximize the
\emph{ideal} projection score.

For \texttt{frame}, the four-term sum $S$ has error below
$2.5\cdot10^{-13}$; its squared modulus has error below
$2.2\cdot10^{-12}$ by \eqref{round:squared-perturbation}, and its modulus
below $2.6\cdot10^{-13}$ by \eqref{round:norm}.  If the stored
$m=\sqrt{4-|S|}$ passes $m>1.25$, the ideal $m>1.24$.
Using $|\sqrt a-\sqrt b|=|a-b|/(\sqrt a+\sqrt b)$ at these lower
bounds gives errors below $3\cdot10^{-13}$ in $m$ and
$\mathrm{top}=\sqrt{4+|S|}$.
In particular $16-|S|^2>6.1$; its error is below
$2.3\cdot10^{-12}$, so its reciprocal has error below
$7\cdot10^{-14}$.  The two pseudoinverse numerators have modulus at
most eight and error below $5.1\cdot10^{-13}$.  For example the
row-times-$\overline S$ product contributes at most the row error
times $|S|$, plus the $S$ error, plus
\eqref{round:complex-multiply}.  Division/scaling then gives
\begin{equation}\label{round:frame-g}
 |\Delta G_{ij}|<10^{-12},\qquad |G_{ij}^{\rm ideal}|<.81,
\end{equation}
the second inequality following from the lower singular bound.

The first constant basis column is exact.  The next column has
numerator error below $1.3\cdot10^{-13}$ and denominator
$\sqrt{4-|S|^2/4}>1.5$ with error below $2\cdot10^{-13}$.
Its ideal entries have modulus at most one, so the quotient error is
below $4\cdot10^{-13}$.  For each of the remaining two columns, use
the actual selected coordinate projection $q$.  In the first such
projection the earlier two products give entry error below
$9\cdot10^{-13}$, and its norm error is below $2\cdot10^{-12}$.
The accepted computed norm is greater than $.449999$, so the ideal
projection is nonzero and normalization gives entry error below
$7\cdot10^{-12}$.  In the last projection the extra column-product
error gives entry error below $1.6\cdot10^{-11}$ and norm error below
$3.5\cdot10^{-11}$.  The same denominator check gives entry error
below $1.2\cdot10^{-10}$.  This argument inductively constructs the
ideal orthonormal complements with the selected coordinates.
The two-product kernel projector has entry error below
$3\cdot10^{-10}$; multiplying two complement entries and $\sqrt6$
gives a W entry error below $8\cdot10^{-10}$.

The following table records the rest of the propagation, with the
source array names.  Its errors include ordinary operation rounding.
\begin{center}
\small
\begin{tabular}{>{\raggedright\arraybackslash}p{.22\linewidth}p{.23\linewidth}p{.45\linewidth}}
\hline
Quantity & Absolute error per entry & Ideal scale used\\
\hline
\texttt{base} & $10^{-11}$ & two factors of modulus at most 2 times G\\
\texttt{fixed} & $2\cdot10^{-9}$ & projector times a factor at most 1.62\\
\texttt{bm}, \texttt{cm} & $2\cdot10^{-12}$ & unit entry times G\\
\texttt{uc.U}, angular derivatives & $2\cdot10^{-13}$ & products of unit phases/trig entries\\
\texttt{K}, \texttt{D} & $4\cdot10^{-9}$ & four W products and base\\
\texttt{mag} & $4\cdot10^{-8}$ & $|D|<4.6$\\
$\sqrt{\texttt{mag}}$ & $5\cdot10^{-9}$ & norm Lipschitz\\
\texttt{jac} & $10^{-8}$ & ideal modulus at most 6\\
flat real J & $10^{-7}$ & dot of D and jac\\
flat c & $5\cdot10^{-8}$ & $(|D|^2-1)/2$\\
pre-equation c & $4\cdot10^{-13}$ & short unit-exponential sums\\
pre-equation coefficients & $6\cdot10^{-14}$ & unit exponentials\\
\hline
\end{tabular}
\end{center}
To verify the larger entries, each base factor $1+c_i\overline a$ has
error below $1.3\cdot10^{-13}$; its multiplication by G and the
two-term sum give the listed base bound.  The larger fixed term uses
the projector error times a factor at most $1.62$, plus a G/unit
combination with error below $3\cdot10^{-12}$, leaving room inside
$2\cdot10^{-9}$.  Four W products cost less than
$4(8\cdot10^{-10})$ plus the unitary and base errors, below
$4\cdot10^{-9}$.  The ideal K entry is at most $\sqrt6$, and the block
operator estimate of Section~\ref{sec:completion} gives
\[
 \|D\|\le\max\{\sqrt6,\|C_0\|A_0/m_B\}<4.6.
\]
Each seed derivative has modulus at most $1.62+.81\cdot4.6<6$;
the individual angular derivatives before the complement/$\sqrt6$
factors have operator norm at most one.  Fixed error, moving-block
error times $.81$, and coefficient error times $4.6$ fit inside
$10^{-8}$.  The flat dot error is then less than
$6(4\cdot10^{-9})+4.6(10^{-8})$ plus a short-dot rounding error,
below $10^{-7}$.  Finally, each $||D_j|-1|\le3.6$; the 16 squared
deviations have total error below $10^{-6}$.  The direct A flat-norm
test adds $10^{-5}$ to its threshold.

\paragraph{The bounds returned by \texttt{get}.}
The stored lower singular bounds subtract $10^{-8}$ and remain
greater than $1.239$.  The A bound adds $10^{-7}$, and L is recomputed
with the inflated top and A bounds before adding another $10^{-7}$.
The seed and angular radii have their explicit $10^{-9}$ and
$10^{-10}$ additions.  They satisfy
\[
 R_s<.203,\quad h<.197,\quad A<2.001,\quad L<6.
\]
With $A_{\max}=\min(2,A+R_s+10^{-8})$ the relevant field scales are
\begin{equation}\label{round:get-scales}
 q<.32,\quad \mathrm{off}<.74,\quad \mathrm{pol}<.068,
 \quad u_{\rm get}<1.26,\quad z<1.5,\quad E<1.
\end{equation}
For example
$q\le .203^2(1+10/1.239^2)+10^{-7}<.32$,
$\mathrm{off}\le .203(\sqrt6+2)/1.239+10^{-7}<.74$, and
$\mathrm{pol}\le .32/(\sqrt6+\sqrt{6-.32})+10^{-7}<.068$.
Then $u_{\rm get}=6h+\mathrm{pol}+10^{-9}$ and the two-dimensional
norm give the z bound.  Substitution of the same bounds in
\[
 E_3=\tfrac12\sqrt{252}\,h^2+\mathrm{pol},\qquad
 E_{\rm phase}=\frac{L+A+1}{2m}R_s^2,
 \qquad e_{\rm off}=E_{\rm phase}+zR_s/m
\]
before their additions gives $E=\|(e_{\rm off},E_3)\|<1$.
All displayed inequalities leave room for the literal/operation errors.

Each field is evaluated from the then-current stored upper/lower
inputs, which are treated as exact in this local step.  The scalar
formulas contain fewer than twenty nonconstant elementary operations,
positive sums/products, divisions by at least $1.239$, or by
$\sqrt6+\sqrt{6-q}>4$, and square roots with arguments at least one
except for norms.  In $6-q$ the relative amplification is less than
$6/5$; norms use \eqref{round:norm}.  Applying
\eqref{round:basic} with these bounds gives error below $10^{-12}$
for each field before its explicit $10^{-9}$, $10^{-8}$, or $10^{-7}$
addition.  The single multiplication error in $R_s^2$ is included.
Positive substitution therefore preserves the upper bounds.
In particular the invalid-$q$ branch cannot occur after a successful
setup at the specified seed radius.

\paragraph{Directional A exclusion.}
\texttt{trial::certificate} rejects a nonfinite direction, then requires
its maximum magnitude to lie between $10^{-100}$ and $10^{100}$.
Division by that maximum produces an exact stored direction with
$|y_i|\le1$.  Its computed constant pairing has error below
$9\cdot10^{-7}$ and each of its 13 derivative pairings has error below
$1.7\cdot10^{-6}$: there are 16 flat rows, four pre-equation rows,
and short dots with ideal J magnitudes at most 28.  Hence the seed and
angular linear supports can together be low by less than
\[
 (.203\sqrt9+.197\cdot4)1.7\cdot10^{-6}
   <2.38\cdot10^{-6}<3\cdot10^{-6}.
\]
The norm $\sqrt{\texttt{sd}}$ approximates
$\|D\circ y_D\|$ within $2\cdot10^{-8}$, by the positive square
bound and the 16 D-entry errors.  All other positive products/norms
and their addition cost less than $10^{-9}$ at these scales.
Thus the actual $.0001$ margin exceeds the whole possible comparison
error.  A failed or skipped proposal does not reject the cell.

\subsection{Root certificates and normal-displacement constants}
\label{round:root}

This subsection verifies the numerical premises of the Root inequalities
in Section~\ref{sec:phase-cover}.  It concerns only successful
\texttt{patch} calls.  The stored center $P$ is exact for these
calculations; \texttt{gen} rejects $|P_j|>4$ or a nonfinite phase.
Each ideal row-edge product is a product of unit exponentials, so its
computed error is below
$2\varepsilon_{\exp}+10^{-15}(1+\varepsilon_{\exp})^2
 <1.3\cdot10^{-13}$.
The six-term edge sum has error below $8\cdot10^{-13}$.
Across the 15 complex edges, division by $2\pi$ gives vector error
less than
\[
 \frac{\sqrt{15}\,8\cdot10^{-13}}{2\pi}<5\cdot10^{-13}.
\]
The squared-modulus, scaling, sum, and final root add less than
$10^{-13}$ at the bounded unit-sum scale.  Thus
$\sqrt{\texttt{val}}$ approximates $\|F(P)\|$ within $10^{-12}$.
The accepted value satisfies $\texttt{val}<10^{-18}$, and the stored
$\mathtt{froot}=\sqrt{\texttt{val}}+10^{-9}$ is an upper bound.

\paragraph{The stored inverse-like matrix M.}
Regard every accepted \texttt{hvec}, \texttt{lvec}, and T entry as an
exact stored number.  Their absolute values are checked to be at most
two.  There are at most 25 good columns.  If $H_V,L_V$ denote these
stored factor arrays, the accumulated M differs from $H_VL_V^*$ by
less than $6\cdot10^{-13}$ per entry: a 25-term dot has absolute
product sum at most 100, giving
$\gamma_{25}100<5.5\cdot10^{-13}$.
There are $36\cdot30$ entries, so the Frobenius difference is below
$2.1\cdot10^{-11}$.

The Gram dots for $H_V$ and $L_V$ have lengths 36 or 30 and absolute
product sums at most 144.  Their errors are below $1.2\cdot10^{-12}$.
After absolute values and at most 25 additions, each Gram row sum
has error below $10^{-10}$: the input error contributes at most
$3\cdot10^{-11}$ and addition error at the 3600 row scale is less
than $2\cdot10^{-11}$.  Consequently the code's
$\mathtt{ah}+10^{-8}$ and $\mathtt{al}+10^{-8}$ bound the largest
eigenvalues of the exact stored-factor Grams.
The product and square root forming
\begin{equation}\label{round:mnorm}
 \mathtt{mnorm}=
  \sqrt{(\mathtt{ah}+10^{-8})(\mathtt{al}+10^{-8})}+10^{-8}
\end{equation}
have error below $5\cdot10^{-12}$, even before the acceptance check
$\mathtt{mnorm}<2$: the root scale is at most 3601, and only a few
positive operations occur.  The final $10^{-8}$ also covers the
$2.1\cdot10^{-11}$ factor-product perturbation.  Hence
\[
 \|M_{\rm stored}\|\le\mathtt{mnorm}<2.
\]
The code's $\mathtt{n0}=\mathtt{mnorm}\,\mathtt{froot}+10^{-9}$
is upward, and so is the first remainder coefficient
$C=\widehat\pi\sqrt8\,\mathtt{mnorm}+10^{-7}$.

\paragraph{Optional tensor improvement.}
For the exact stored M, let $A_*$ be the $30$-by-$30$ row Gram of
the tensor map $\mathcal S$ in Section~\ref{sec:phase-cover}.
Its $(a,b)$ entry is an M-column dot product times
\[
 (\mathbf e_i-\mathbf e_j)\mathbin{\cdot}
 (\mathbf e_{i'}-\mathbf e_{j'}),
\]
with the edge sign repeated for the two real edge parts.  The sign
factor is an integer of magnitude at most two.  A length-36 dot with
products bounded by four therefore has error below
$3\cdot10^{-12}$ after that multiplication.  Since $\|M\|<2$,
the ideal Gram entries have magnitude at most eight.

The Jacobi output U is used only if all its stored entries have
absolute value at most two.  For that exact stored U, the residual
Gram $U^*U-I$ has computed entry error below $10^{-12}$, since its
dots have length 30 and absolute sum at most 121 including the
identity term.  Its 900 entry errors have Euclidean total at most
$3\cdot10^{-11}$.  The residual Frobenius norm is at most 3631 at
these loose entry bounds; its accumulation/root error is below
$8\cdot10^{-10}$.  Thus
\[
 \mathtt{beta}=\operatorname{fl}(\|U^*U-I\|_F)+10^{-8}
\]
is an upper bound for the operator defect.

The computed product X differs from $A_*U$ by less than
$2\cdot10^{-10}$ per entry: input perturbation contributes at most
$30\cdot2\cdot3\cdot10^{-12}$, and its dot has absolute product
sum at most 481.  Exact product entries are at most 480.  The further
product $U^*X$ differs from $U^*A_*U$ by less than
$1.3\cdot10^{-8}$ per entry: at most
$30\cdot2\cdot2\cdot10^{-10}$ comes from X, and the length-30
dot contributes less than $2\cdot10^{-10}$.
Each absolute row sum therefore has error below $10^{-6}$,
including addition rounding at row scale below $900000$.

If $\mathtt{beta}<.1$, the exact stored U is invertible and
$\|U^{-1}\|^2\le(1-\mathtt{beta})^{-1}$.
The exact matrix $U^*A_*U$ is symmetric.  The inflated computed
absolute-row bound $\mathtt{row}+.0002$ therefore gives
\[
 \|A_*\|\le
 \frac{\mathtt{row}+.0002}{1-\mathtt{beta}}.
\]
The update
\begin{equation}\label{round:tensor-c}
 C\leftarrow\min\left\{C,
  \widehat\pi\sqrt{\frac{6(\mathtt{row}+.0002)}
                              {1-\mathtt{beta}}}+.00001\right\}
\end{equation}
is consequently safe.  The denominator exceeds $.9$; at the displayed
row scale the scalar evaluation error is below $10^{-8}$, smaller
than its final $.00001$.  No eigensolver accuracy has been assumed.

\paragraph{The horizontal-projector identity.}
An entry of the exact double-centering projector is the product of
the two appropriate centering-matrix entries.  The computed
\texttt{PI} differs by less than $10^{-15}$ per entry.  In the
matrix $L_P$, each row has only 12 nonzero entries, with error below
$1.3\cdot10^{-13}$ from the edge products.  Forming
\texttt{LP} gives error below $5\cdot10^{-12}$ per entry against
$L_P\Pi$; the ideal entry magnitude is at most $8.4$.
In the residual
\[
 \Pi-TT^*-M L_P\Pi
\]
each computed entry then has error below $10^{-9}$.
The propagated LP error is at most
$30\cdot2\cdot5\cdot10^{-12}=3\cdot10^{-10}$;
the remaining PI, T-product, and short accumulation errors fit in
the remainder of the budget, with absolute sum below 560.
The 1296 entry errors have Frobenius total below $3.6\cdot10^{-8}$.
Norm accumulation is below $10^{-8}$ even with residual norm 20161.
Thus the actual $2\cdot10^{-6}$ addition in \texttt{id} suffices.

For $T^*T-I$, length-36 dots have errors below
$1.2\cdot10^{-12}$.  With $k\le10$, their aggregate and norm error
is below $10^{-10}$.  Projecting a stored T column by the row/column
mean formula has component error below $10^{-13}$: its row, column,
and grand sums have lengths 6, 6, and 36 and entries bounded by two.
The subtraction and Frobenius calculation of
$(I-\Pi)T$ have error below $10^{-11}$.
The $10^{-8}$ additions in \texttt{ort} and \texttt{hor} cover these.
We have therefore established, on every successful Root path,
\begin{align}
 \|F(P)\|&\le\mathtt{froot},&
 \|M\|&\le\mathtt{mnorm},&
 \mathtt{mnorm}\mathtt{froot}&\le\mathtt{n0},\nonumber\\
 \|\Pi-TT^*-ML_P\Pi\|&\le\mathtt{id},&
 \|T^*T-I\|&\le\mathtt{ort},&
 \|(I-\Pi)T\|&\le\mathtt{hor}.\label{round:root-invariants}
\end{align}
The acceptance thresholds are
$\mathtt{id}<10^{-4}$,
$\mathtt{ort},\mathtt{hor}<10^{-5}$, $C<20$, and
$\mathtt{mnorm}<2$.  In particular
\begin{equation}\label{round:t-operator}
 \|T\|\le\sqrt{1+\mathtt{ort}}<1.00001.
\end{equation}
Signed permutation transport of stored P and T is exact before
dephasing.  The transported inequalities of
Section~\ref{sec:phase-cover} therefore use the same constants.

\subsection{Integer lifts, dephasing, and sphere distances}
\label{round:dephasing}

For $|p_j|\le4$, the three operations in unwrapped dephasing
$p_{ij}-p_{i0}-p_{0j}+p_{00}$ have entry error below
$8\cdot10^{-15}$; their intermediate magnitudes are at most 8, 12,
and 16.  Subtracting the integer selected by \texttt{nearbyint}
gives error below $10^{-14}$ against the exact dephase minus that
same integer.  A wrapped difference of two such arrays has error
below $3\cdot10^{-14}$ per entry against the ideal difference with
the combined selected integers.  The one-sided Candidate version,
whose other array is already a stored dephased phase, has the same
bound.  These comparisons do not assume identical rounding choices
for an independently formed ideal phase.

For a stored wrapped array, $|z_j|\le1/2$.  Its row and column sums
have error below $4\cdot10^{-15}$, and its length-36 grand sum below
$1.4\cdot10^{-13}$.  Dividing by 6 or 36 and combining the means
gives projection-operation error below $10^{-14}$ per component.
Projection contracts the prior input error.  Consequently the
computed projected difference has Frobenius error below $10^{-11}$.
All distance-square and norm errors at the successful cutoffs below
are negligible compared with $10^{-7}$.

Thus \texttt{Atlas::find}'s returned $\sqrt s+10^{-7}$ is an upper
distance for the selected ideal projected lift.  A merge uses radius
$r+\text{distance}+10^{-8}$ and checks a cap with $2\cdot10^{-8}$
slack.  Its sum-rounding error is below $10^{-14}$.
The same arithmetic applies to \texttt{assign} and to the local
sphere-radius additions.  Searching only some buckets or transforms
may miss a merge; it cannot certify a merge that fails these tests.

\subsection{Candidate conversion and ellipsoidal coverage}
\label{round:candidate}

The coverage estimates here are needed only after successful
Candidate preparation, at final unitary-grid levels $N_u\ge64$.
Let $m_j=|D_j|$ denote an ideal comparison-block magnitude.
The stored magnitude test $\sqrt{\texttt{mag}_j}>.15$ implies
$m_j>.14999998$.  Normalizing D by its computed magnitude has
unit-direction error below
\[
 \frac{2\cdot4\cdot10^{-9}}{.14999998}
      +\text{norm/division rounding}<5.5\cdot10^{-8}.
\]
Candidate additionally checks a computed chord smaller than
$10^{-7}$ between this normalized value and the exponential of its
proposed phase.  With \eqref{round:exp-budget}, the true unit chord
error is less than $1.56\cdot10^{-7}$.
If a circular phase distance in cycles is $v\in[0,1/2]$, its chord
is $2\sin(\pi v)\ge4v$.  Across 16 entries the phase-center
correction therefore has norm below $1.56\cdot10^{-7}$, and in
particular below the $2\cdot10^{-6}$ allowance used in
Section~\ref{sec:phase-cover}.  The nine seed entries use exact stored
seed centers and an exact integer subtraction.

For the D-position coefficients X, the normalized-direction error
times the Jacobian bound six, plus the $10^{-8}$ Jacobian error and
short-dot rounding, is below $10^{-6}$ per entry.  The other used
X entries are exact.  The $\mathtt{mi}$ subtraction of $10^{-6}$
therefore gives a positive magnitude lower bound greater than $.148$,
and
$\mathtt{invmag}_j=1/(\widehat m_j-10^{-6})+10^{-9}$
is an upper bound for $1/m_j$, below $6.8$.

The flat-sum error is below $10^{-6}$, so the $10^{-5}$ addition in
\texttt{slack}, followed by its additional $10^{-10}$, gives the
required nonnegative weighted chord-square upper bound.  In the
passed branch $q<3.8$ and $1-q/4>.049$.  The arithmetic in dividing
\texttt{slack} by \texttt{mi} has error below $10^{-12}$:
\texttt{slack}$<4.1$ and the denominator is above $.148$.
The ratio $(1-q/4)^{-1/2}$ has error below $10^{-12}$ before its
$10^{-7}$ addition; the denominator lower bound bounds the
derivative, and $q/4$ is a binary scaling.  The remaining positive
products/roots, and the division by $(2\pi)^2$ in \texttt{d2}, each
have error below $10^{-12}$ before their $10^{-8}$ additions.
In particular the stored quantities are upper bounds with
\begin{equation}\label{round:candidate-scales}
 \mathtt{arc}<8.8,\qquad \mathtt{weighted}<10,
 \qquad \mathtt{invmag}_j<6.8.
\end{equation}

\paragraph{The direct distance bound.}
If \texttt{nd}$\le.145^2$, the ideal horizontal difference $\Delta$
has norm below $.146$.  Its computed squared norm has error below
$10^{-10}$ by Section~\ref{round:dephasing}.
Each pairing of $\Delta$ with an X column has error below
$7\cdot10^{-7}$: the D-part column error has norm at most
$4\cdot10^{-6}$, the ideal column norm is below 25, and
\[
 .146\cdot4\cdot10^{-6}+25\cdot10^{-11}
      +\text{dot rounding}<7\cdot10^{-7}.
\]
At $R_s<.203$ and $h<.05$, the two linear support terms have combined
error below $10^{-6}$.  The D-subnorm error is below
$2\cdot10^{-11}$ and its infinity-norm error below $10^{-11}$;
the latter remains harmless even after multiplication by
$8.8^3/6$.  For the alternative support, the weighted norm on exact
stored positive \texttt{invmag} has component perturbation at most
$\sqrt{6.8}\,10^{-11}$, plus positive square/norm rounding.  The
weights themselves are upper bounds.
Thus either support can be low by less than $2\cdot10^{-6}$, and
their minimum has the same property.

The code adds $.0001$ to that minimum before division by $\pi$.
This covers the support error, the $10^{-10}$ squared-distance error,
the error in division by $\pi$, and the remaining positive additions
in the square-root argument.  The final $.00001$ covers the center
correction $2\cdot10^{-6}$ and square-root rounding.  Consequently
\texttt{s0} is a valid upper bound for the true aligned solution
distance.

\paragraph{Tangent offsets and external errors.}
The computed \texttt{off} approximates $2\pi T^*\Delta$ within
$10^{-9}$ in norm, by \eqref{round:t-operator}, the projected
difference error, and short-dot rounding.  The error in the whole
X array is at most $\sqrt{208}\,10^{-6}$ in Frobenius norm; there are
13 columns and 16 affected positions.  Multiplication by $T^*$ and
the less than $2\cdot10^{-13}$ rounding of each dot give
\begin{equation}\label{round:ab-error}
 \|\mathtt{AB}-T^*X_{\rm ideal}\|<1.5\cdot10^{-5}.
\end{equation}
The parameter-vector norm is below
$\sqrt{.203^2+4(\pi/64)^2}<.226$, so substitution of stored AB
costs less than $3.4\cdot10^{-6}$ radians.  Formation of the corner
offsets adds less than $10^{-11}$, using coefficients of magnitude
below 25.  The same bound applies to convex combinations of these
corners.

Using raw rather than projected angular displacement costs at most
\[
 \mathtt{hor}\sqrt{R_s^2+\mathtt{arc}^2}<8.81\cdot10^{-5}
 \quad\hbox{radians},
\]
and the center correction costs at most
$2\pi\sqrt{1+\mathtt{ort}}\,2\cdot10^{-6}<1.27\cdot10^{-5}$
radians.  The row-subnorm factor multiplying the sine remainder has
downward arithmetic error below $10^{-10}$ even after its
$\mathtt{arc}^3/6$ factor, by positive square/norm propagation.
Adding these errors and \eqref{round:ab-error}'s effect gives less
than $.00012$ radians in either appropriate tangent bound, excluding
the covariance/ell error treated next.  The code adds $.00015$
radians and then further cycle-unit slack of $.00002$.

\paragraph{Covariance domination.}
For the first covariance, compare the exact stored symmetric G with
the mathematical covariance using the \emph{stored} AB, $R_s$, and E,
and rational $\beta=.45$.  Each nine-term AB outer-product dot has
absolute sum at most $9\cdot25^2=5625$ and error below
$1.3\cdot10^{-11}$.  Its multiplier $1.45R_s^2$ is below $.06$.
The D2 dots have 16 products, bounded even by four each.
Including coefficient translation, products, sums, and the intended
diagonal addition gives total entry error below $10^{-11}$.

For the second covariance, use stored \texttt{weighted} and
\texttt{invmag}, with rational $\beta=.8$.  Each Dweighted dot has
absolute sum less than
$16(1.00001)^2\cdot6.8<110$.
Its two-product-per-term accumulation has error below
$5\cdot10^{-13}$ by \eqref{round:dot} with one extra factor.
The multiplying coefficient is below 225.  Coefficient translation
and the subsequent arithmetic, even at absolute intermediate scale
26000, give total entry error below $10^{-9}$.
Both G arrays are explicitly stored symmetrically.  The added
$10^{-7}I$ dominates $k$ times either entry-error bound for $k\le10$.
Hence the exact stored G is positive semidefinite and dominates the
required covariance in each case.  This conclusion does not depend
on numerically observed positive eigenvalues.

\subsection{The spectral and inverse-quadratic-form bounds}
\label{round:ell}

\begin{lemma}[The stored spectral upper bound]\label{round:spectral}
For an exact stored symmetric $k$-by-$k$ matrix G with $k\le10$ and
$|G_{ij}|\le10$, \texttt{spectralUpper} returns an upper bound for
$\|G\|$.
\end{lemma}

\begin{proof}
An absolute row sum has error below $3\cdot10^{-13}$ by
\eqref{round:dot} or direct positive-sum propagation.  The code adds
$10^{-10}$ to their maximum.
For its other bound, put $m=\max_{ij}|G_{ij}|$ and $A=\|G^2\|_F$.
Each computed G-squared dot has error below
$3\cdot10^{-14}m^2+10^{-297}$: it has ten products and absolute sum
at most $10m^2$.
Symmetry implies $A\ge m^2$, because one diagonal sum of squares
contains a term $m^2$.  The 100 dot errors and the norm evaluation
therefore give
\[
 |\operatorname{fl}(\sqrt{\mathtt{fro4}})-A|
       <4\cdot10^{-13}A+10^{-139}.
\]
This remains valid when individual off-diagonal dots cancel or
underflow.  Taking the next square root, by the inequality used in
Lemma~\ref{round:accumulation}, costs less than
$4\cdot10^{-13}\sqrt A+10^{-69}$ before its own tiny rounding.
Since $\sqrt A\le100$, the total possible downward error is less
than $10^{-10}$.  The code adds $2\cdot10^{-7}$.
The exact fourth-norm quantity $\sqrt{\|G^2\|_F}$ bounds $\|G\|$.
Both returned candidates are upper bounds, so their minimum is also
an upper bound.  In particular its value \texttt{gu} is nonnegative
and at most $100.1$.
\end{proof}

The \texttt{ell} guard checks symmetry and $|G_{ij}|\le10$, and checks
each offset component to have magnitude at most ten.  Its invalid
return value 100 cannot certify a cell: after conversion to cycles
it exceeds 10, so the associated bootstrap path returns its
nonsuccess sentinel.  A different valid bound may still succeed.
For valid data the initial triangle bound is
$\sqrt{\mathtt{mx}}+\sqrt{\mathtt{gu}}+10^{-8}$.
Here $\mathtt{mx}$ is the maximum squared offset norm.  Its norm
error is below $10^{-13}$ and the triangle bound is upward and
smaller than 43.

For each proposed resolvent parameter, interpret stored
$\mathtt{gam}=\gamma$ as exact.  The proposal formula and the input
scales give $0<\gamma<418$.  Since \texttt{gu} is an upper spectral
bound,
\begin{equation}\label{round:gap}
 \mathtt{gap}=\operatorname{fl}(\gamma-\mathtt{gu}-10^{-9})
       \le\gamma-\lambda_{\max}(G).
\end{equation}
The two subtraction errors are below $2\cdot10^{-13}$, absorbed by
$10^{-9}$.  Only a stored gap of at least $10^{-6}$ is used.
Allowing literal translation still gives a gap above
$.999\cdot10^{-6}$ for scale estimates.

For one offset vector U, the trial solve is used only if each stored
$|z_i|<10^8$.  Let $s_z=\max(1,\|z\|_\infty)$.
Its exact residual is $r=U-\gamma z+Gz$.  The following local bounds
hold:
\begin{center}
\small
\begin{tabular}{p{.25\linewidth}p{.36\linewidth}p{.29\linewidth}}
\hline
Computed quantity & Error against exact stored-input expression & Scale\\
\hline
each residual component & $2\cdot10^{-12}s_z$ & absolute sum $\le528s_z$\\
residual norm & $2\cdot10^{-11}s_z$ & norm $\le\sqrt{10}\,528s_z$\\
U--z dot & $3\cdot10^{-13}s_z$ & absolute sum $\le100s_z$\\
U norm & $10^{-13}$ & norm $\le\sqrt{1000}$\\
\hline
\end{tabular}
\end{center}
For example the component bound follows from at most twelve terms
and the scale 528 using \eqref{round:dot}; its aggregate contributes
at most $\sqrt{10}\,2\cdot10^{-12}s_z$ to the norm, and
\eqref{round:norm} adds less than $5\cdot10^{-12}s_z$.
The dot and U-norm bounds are the same lemmas at lengths ten.

The exact inverse quadratic form satisfies
\begin{equation}\label{round:resolvent-residual}
 U^*(\gamma I-G)^{-1}U
 \le U^*z+\frac{\|U\|\|r\|}{\gamma-\lambda_{\max}(G)}.
\end{equation}
Indeed $(\gamma I-G)^{-1}U=z+(\gamma I-G)^{-1}r$.
The code forms the right-hand-side upper using
\begin{equation}\label{round:ell-upper-code}
 \big(\mathtt{dotv}+10^{-10}s_z\big)
 +\frac{(\sqrt{\mathtt{un}}+10^{-10})
         (\sqrt{\mathtt{res}}+10^{-10}s_z)}{\mathtt{gap}},
\end{equation}
with the shown binary64 operations.
The rounded dot part exceeds $U^*z$ by at least
$9\cdot10^{-11}s_z$ and has magnitude at most $101s_z$.
The first rounded norm factor is at least
$(1+2\cdot10^{-12})\|U\|$, because $\|U\|\le\sqrt{1000}$;
the second is at least $\|r\|$.
Together with \eqref{round:gap}, product/division rounding therefore
leaves a quotient B such that the required exact quotient in
\eqref{round:resolvent-residual} is at most $(1-10^{-12})B$.
This quotient is finite and below $6\cdot10^{18}$ by the displayed
norm scales and the gap lower bound.  The final addition's possible
error is at most $u(101s_z+B)+\tau$, which is smaller than the
surplus $9\cdot10^{-11}s_z+10^{-12}B$.
Thus \eqref{round:ell-upper-code} is an upper bound even when its
dot part is negative and cancellation occurs.

Taking the maximum with zero and over the offsets gives
\texttt{high}.  The ellipsoid formula of
Section~\ref{sec:phase-cover} then uses
$\sqrt{\gamma(1+\mathtt{high})}+10^{-6}$.
There is no overflow at the stated scales.  If this candidate can
lower the existing bound below 43, its exact unrounded root is less
than 44; its scalar downward error is then below $10^{-12}$,
covered by $10^{-6}$.  A candidate that cannot improve the existing
upper need not have a tighter absolute estimate.  The routine adds
another $10^{-6}$ at return.  Cholesky accuracy has not entered this
argument; it merely proposes the checked stored z.

\paragraph{Bootstrap and saved radii.}
Together with Section~\ref{round:candidate}, this establishes the
tangent upper whenever its bootstrap path is usable.
The bootstrap guard requires $0\le\mathtt{init}<3$, $0\le t<10$,
and $C<20$, with previously certified nonnegative Root constants.
It first adds $.00002$ to t and $10^{-6}$ to the distance upper.
At a current upper $s\le3.00001$, the normal polynomial
$\mathtt{n0}+\mathtt{id}s+Cs^2$ is below 181.
Its positive-operation error, the tangent-factor error, and the
two-component norm error total less than $10^{-12}$; there are only
a few products/sums, and \eqref{round:norm} applies at norm scale 181.
The update's $10^{-8}$ addition is sufficient.  Accepting an
improvement therefore preserves the upper inductively, and stopping
keeps the preceding valid upper.  This is a finite sequence of up
to 30 upper-bound improvements, not a convergence assertion.
The subsequent positive additions in saved sphere radii preserve
the same inequalities.

\subsection{Refinement relations, tangent grids, and radius updates}
\label{round:refinement}

The B inputs are genuine earlier Root certificates.  In addition,
\texttt{validate} enforces $0\le R\le.126$, $0\le C<20$,
$0\le\mathtt{n0},\mathtt{froot}<10^{-6}$,
$0\le\mathtt{id}<10^{-4}$,
$0\le\mathtt{ort},\mathtt{hor}<10^{-5}$,
$0\le\mathtt{mnorm}<2$, $0\le k\le10$, and $|p_j|\le4$.
The stored T entries pass the range check $|T_{ja}|\le1.01$;
their stronger operator bound is
\eqref{round:t-operator}, supplied by the genuine Root certificate.

\paragraph{The relation initialization.}
On a passed raw-difference/cut path,
Section~\ref{round:dephasing} gives a projected-difference error below
$10^{-11}$ in norm, and the computed cutoff implies
$\|d_q\|<.251$ for the ideal chosen lift.
More precisely, \texttt{cut} is at most the stored \texttt{.249}.
The passed norm check, with its \texttt{1e-6} allowance, therefore
gives $\widehat{\|d_q\|}<.249002$, including literal and operation errors.
Each computed $\eta_{q,a}=T_a^*d_q$ has error below
$1.1\cdot10^{-11}$: multiply the projection error by the
norm-$<1.00001$ T column and include a length-36 dot.
For $n_q=\|(I-TT^*)d_q\|$, combine the $d_q$ error with
\[
 \|T\|\sqrt{10}\,1.1\cdot10^{-11}
\]
and the short entrywise products and final norm.  The result is below
$5\cdot10^{-11}$.  The actual $10^{-8}$ addition to
\texttt{nq} is sufficient.

Writing $o=\mathtt{ort}_P$, the ideal quantities in
Section~\ref{sec:refinement} are bounded by the stored fields
\[
 \begin{split}
 \mathtt{total}&=R+\widehat{\|d_q\|}+10^{-8},\\
 \mathtt{nu0}&=\mathtt{n0}_P+\mathtt{id}_P R+C_P R^2
                    +\mathtt{nq}+10^{-8},\\
 \mathtt{lam}&=\mathtt{id}_P+C_P\mathtt{total}+10^{-8},\\
 \mathtt{e}&=\mathtt{mnorm}_P\mathtt{froot}_Q+10^{-7}.
 \end{split}
\]
Their ordinary positive-operation errors are below $10^{-13}$,
and their input errors have already been covered.
The following loose scales hold on every passing path:
\begin{equation}\label{round:b-init-scales}
 \mathtt{total}<.377,\qquad \mathtt{nu0}<.58,\qquad
 \mathtt{lam}<7.55,\qquad \mathtt{e}<2.2\cdot10^{-6}.
\end{equation}
For the second bound, use
$\|I-TT^*\|\le1+o$ and $\|d_q\|<.251$; the parent normal polynomial
at $R\le.126$ is below $.318$.
For the first, the sharper stored-norm bound above gives
\[
 \mathtt{total}<.126+.249002+10^{-8}+10^{-12}<.375003<.377,
\]
where $10^{-12}$ covers literal translation and the two additions.
The remaining inequalities follow from these range bounds.

For the cross-normal factor $\sigma$, the code first projects each
child T column, giving S, and then forms
$W=(I-TT^*)S$.  A projected component has error below $10^{-13}$.
Its dot with a parent column has error below $7\cdot10^{-13}$
against the exact projected dot: the projection error has norm
below $6\cdot10^{-13}$, and the remaining dot rounding is small.
At most ten such terms in W give component error below $10^{-11}$.
Exact W column norms are below $1.01$, since both projection and
the checked T operators are bounded.  Thus the W-column Gram has
entry errors below $2\cdot10^{-10}$: two column perturbations each
of norm at most $6\cdot10^{-11}$, times $1.01$, plus a length-36
dot suffice.  The spectral perturbation is below $2\cdot10^{-9}$.
The Gram is explicitly stored symmetrically and passes
$|G_{ij}|\le2$, so Lemma~\ref{round:spectral} applies.
The returned
\[
 \mathtt{sig}=\sqrt{\mathtt{spectralUpper}(G)+10^{-7}}+10^{-8}
\]
is an upper bound.  The inner addition exceeds the Gram perturbation
and any final scalar addition rounding; the outer one exceeds the
square-root rounding.  Its value is below $4.6$, already from the
absolute-row bound.  For zero child tangent columns, $\sigma=0$
is exact.

\paragraph{The upper-bound iteration and the rho search.}
Interpret a passed input $t$ as an exact stored upper bound on the
desired tangent distance, or as a threshold for the entire interval
of distances below it.  The guard is $0\le t\le2$.
The initial increase of $t$ by $2\cdot10^{-8}$ and the
$10^{-9}$ addition in
$z=\sqrt{1+o}\,t+10^{-9}$ ensure the parent tangent upper.
Both initial distance choices, \texttt{total} plus $10^{-8}$ and
$\operatorname{hp}(z,\mathtt{nu0})+10^{-8}$, are valid.
With $o_c=\mathtt{ort}_Q$, each iteration forms upper bounds
corresponding to
\[
 \begin{split}
 v&=
 \min\{\mathtt{nu0},(1+o)d,\mathtt{lam}\,d+\mathtt{e}\}
                     +10^{-8},\\
 e_p&=(1+o)\sqrt{(1+o)(1+o_c)}\,t+\mathtt{sig}\,v+10^{-8},\\
 d_1&=\operatorname{hp}(z,v)+10^{-8},\\
 d_2&=\operatorname{hp}\bigl(\sqrt{1+o_c}\,e_p,\,
              \mathtt{n0}_Q+\mathtt{id}_Qd+C_Qd^2\bigr)+10^{-8}.
 \end{split}
\]
Here v is a computed upper for the true normal distance.
The distance update takes the minimum of the last two bounds only
if it improves the previous upper.
Throughout, the stored current distance is below $.378$,
the normal upper below $.39$, $e_p<3.9$, and the child normal
polynomial below $2.9$.
The $(1+o)d$ choice alone supplies the normal scale.
At these scales each positive product/sum, factor, and
two-component norm has total local downward error below $10^{-13}$,
including the child polynomial's error before its use in the norm.
This follows from fewer than twenty elementary operations per field,
the displayed magnitudes, and \eqref{round:norm}; no small norm
appears in a denominator.  Thus every $10^{-8}$ addition, and the
$10^{-9}$ in z, preserves the upper.  The final $10^{-8}$ addition
also does so.

If the rho search returns $\rho>10^{-7}$, a positive lower-endpoint
update must have occurred.  Its final stored \texttt{lo} came from
an actual trial with
$\mathtt{eval}(\mathtt{lo})<\mathtt{target}-2\cdot10^{-7}$.
Every actual tangent distance at most this \texttt{lo} is covered
by that upper computation.  The returned
$\rho=\max(0,\mathtt{lo}-10^{-8})$ is at most \texttt{lo};
subtraction rounding at these scales is far smaller than $10^{-8}$.
This argument requires no monotonicity of rounded evaluations and
no assertion that bisection finds a maximal valid threshold.
For child-radius growth, \texttt{eval} is directly recomputed at
the assigned tangent-distance upper, inflated, and checked against
the target.

\paragraph{The base grid and the outside test.}
Use the stored positive width $w_B$ as an exact grid parameter.
The allowed range is $10^{-4}\le w_B\le.05$.
The stored grid radius
$E=\sqrt{1+\mathtt{ort}_P}\,R+10^{-7}$ exceeds the true tangent
radius by more than $.999\cdot10^{-7}$.
For $L=\lceil E/w_B+.500001\rceil$, the quotient scale is at most
1261 and its error is below $10^{-12}$.
The $.000001$ excess over the half-cell term is much larger,
so the integer base range covers $[-E,E]$ in every coordinate.
In particular $L\le1262$.
Exact base centers have magnitude below $.203$ per coordinate.
Their computed products have error below $10^{-16}$.
Halving a stored width is exact here; after at most six subdivision
levels, each computed center differs from the intended exact dyadic
child center by less than $10^{-15}$ per coordinate.

For a box center c and halfwidth h, its true minimum coordinate
distance from zero is $(|c_a|-h)_+$.
The code instead forms
$(|\widehat c_a|-h-10^{-8})_+$, whose exact value is no larger than
the true one after all center and elementary-operation errors.
If a box intersects the true tangent ball of radius $E_0$,
the exact sum of these squared lower distances is at most $E_0^2$.
The square-sum rounding is below $4\cdot10^{-17}$ at
$E_0\le.127$, while
\[
 E^2-E_0^2
 \ge 2E_0(.999\cdot10^{-7})+(.999\cdot10^{-7})^2
 >9.9\cdot10^{-15}.
\]
This also handles $E_0=0$.
Rounding of $E^2$ is smaller still.  Consequently \texttt{outside}
cannot reject an intersecting box.

\paragraph{Farthest corners and marking.}
The corner-distance norm before its final margin has error below
$5\cdot10^{-11}$ against the exact box and exact $\eta_q$:
use the $\eta_q$ component errors, center errors, and
\eqref{round:norm}.  Its $10^{-8}$ addition is sufficient.
For \texttt{markRec}, each coordinate radius
\[
 |\widehat c_a-\widehat\eta_{q,a}|+w_B/2+10^{-8}
\]
already exceeds the true farthest coordinate distance.
The accumulated square sum s, followed by
$\sqrt s+2\cdot10^{-8}$, therefore bounds the true farthest tangent
distance, including its norm-rounding error.
For $k>0$, the final accumulated s has also passed comparison with
$(\rho-5\cdot10^{-8})^2$ and $\rho\ge5\cdot10^{-8}$.
At $\rho<.13$, square/comparison/root rounding costs less than
$10^{-14}$ in distance; hence
\[
 \sqrt s+2\cdot10^{-8}
 \le \rho-3\cdot10^{-8}+10^{-14}<\rho.
\]
For $k=0$, the marked distance is $2\cdot10^{-8}$, which is below
the required positive rho threshold.
Thus a marked box is wholly covered, and the child-radius update
receives a valid tangent upper within rho.
Heuristic integer ranges can miss marks but cannot mark a box
without this whole-box test.
Split-box use of the separate corner test has an additional
$10^{-8}$ comparison slack, and a direct successful snap assignment
uses its own inflated \texttt{eval} return and target comparison.
The triangle merges following assignment were certified in
Section~\ref{round:dephasing}.

\subsection{Single-vector filters and vector-ball construction}
\label{round:single}

The C stage validates its Root and $0\le r\le.03$.
At every bin resolution $M\ge6$, the halfwidth, Euclidean bin
radius X, one-norm bound, and infinity-norm bound are formed as
\[
 \begin{aligned}
 a&=.5/M+10^{-12},&
 X&=\sqrt{5-1/6}\,a+10^{-12},\\
 x_1&=5a+10^{-12},&
 x_\infty&=1.5a+10^{-12}.
 \end{aligned}
\]
Their operations have error below $10^{-14}$, and their
$10^{-12}$ additions are upward.
They satisfy $X<.19$, $x_1<.42$, and $x_\infty<.126$.
The stored tangent and normal bounds
\[
 r_t=\sqrt{1+\mathtt{ort}}\,r+10^{-9},\qquad
 r_n=\mathtt{n0}+\mathtt{id}r+Cr^2+10^{-9}
\]
are upper bounds, with $r_t<.030001$ and $r_n<.019$.
The square-root/product and positive-polynomial errors are below
$10^{-13}$.
The constant $K=1.00001$, even allowing its literal translation
error, exceeds the needed operator bound: the Root residual gives
\[
 \|e(P)/\sqrt6\|
 \le\sqrt{1+(2\pi\sqrt2/6)\,10^{-6}}
 <1.00000075<K.
\]

For bin or ball centers, put, as in
Section~\ref{sec:exclusion},
\[
 q_{ji}=e(u_j-P_{ji})/\sqrt6,\qquad S_i=\sum_jq_{ji},\qquad
 c_i=(|S_i|^2-1)/(4\pi).
\]
Let $J_{ij}=-\Im(\overline S_iq_{ji})$ and
$D_{ij}^{\rm re}=\Re(\overline S_iq_{ji})$.
The following errors are uniform for the called arrays:
\begin{center}
\begin{tabular}{lll}
\hline
Quantity & Entry error & Ideal magnitude bound\\
\hline
q & $10^{-13}$ & $1/\sqrt6$\\
S & $7\cdot10^{-13}$ & $\sqrt6$\\
$|S|^2$ & $4\cdot10^{-12}$ & 6\\
$|S|$ & $10^{-12}$ & $\sqrt6$\\
c, J, $D^{\rm re}$ & $10^{-12}$ & J and $D^{\rm re}$ at most 1\\
B & $7\cdot10^{-12}$ & at most $\sqrt6\,1.00001$\\
\hline
\end{tabular}
\end{center}
Indeed the scaled exponential error is actually less than
$2.5\cdot10^{-14}$ before the loose $10^{-13}$ allowance.
Six terms and addition rounding give the S bound.
Equation~\eqref{round:squared-perturbation} gives less than
$2\sqrt6(7\cdot10^{-13})$ plus tiny arithmetic, below
$4\cdot10^{-12}$.
Norm Lipschitz gives the magnitude bound.
The companion products cost at most
$(7\cdot10^{-13})/\sqrt6+\sqrt6\,10^{-13}$
plus product rounding, below $10^{-12}$.
Finally $B_{ia}=\sum_jJ_{ij}T_{ji,a}$, where $T_{ji,a}$ is the
entry at phase position $(j,i)$ in tangent column a, has six terms, each T
entry bounded by $1.00001$, so its input and short-dot errors are
below $7\cdot10^{-12}$.

\paragraph{Direct bin filters.}
The scalar magnitude-change bound
$2\pi(x_1/\sqrt6+r)$ and its squared thresholds have arithmetic
error below $10^{-10}$ at the displayed scales.
The same holds for the aggregate norm threshold and the accumulated
distance of $(|S_i|-1)_i$; the six magnitude errors, multiplied by
bounded factors in their squares, are below $10^{-10}$ in total.
The lower scalar test is used only when the computed magnitude-change
bound is below one.
If the required ideal bound crosses one solely because of rounding,
the tested lower square is of order at most $(10^{-10})^2$,
and the code subtracts $2\cdot10^{-7}$ from it.
Otherwise the squared lower threshold can be propagated on the
correct side of one with the same $10^{-10}$ allowance.
Thus this conditional lower test is safe also at the crossing.

The scalar linear-support sums use at most five absolute J
coefficients and a six-component J norm.
Their input errors, the two alternative Taylor remainders, and
the error in the maximum S magnitude give total possible
underestimation below $10^{-9}$.
For example their linear input effects are at most
$5\cdot10^{-12}a+\sqrt6\,10^{-12}r$; the remainder factors
use only $X<.19$, $x_\infty<.126$, $r\le.03$, and
$|S_i|\le\sqrt6$, and hence multiply the S/J perturbations by
bounded factors below ten.
The remaining positive scalar arithmetic is below $10^{-10}$.
All these direct exclusions use a $2\cdot10^{-7}$ comparison margin.

\paragraph{The directional single-vector support.}
The heuristic solve is accepted only with finite coordinates and a
maximum magnitude between $10^{-100}$ and $10^{100}$.
After division by that maximum, regard the resulting stored
$|y_i|\le1$ as the exact test direction.
Define the exact stored-coefficient expressions
\[
 \begin{aligned}
 a_y&=\sum_i y_i c_i,& h_a&=\sum_i y_iB_{ia},\\
 z_j&=\sum_i y_iJ_{ij},&
 \widetilde z_j&=\sum_i y_iD_{ij}^{\rm re},\\
 N_J&=\sum_{i,j}y_i^2J_{ij}^2,& N_S&=\sum_i y_i^2|S_i|^2,\\
 M_S&=\max_i y_i^2|S_i|^2,& m_y&=\max_i|y_i|,\\
 M_z&=\max_j(z_j^2+\widetilde z_j^2).
 \end{aligned}
\]
The support compared with $|a_y|$ is the stored evaluation of
\begin{align}
 &a\sum_{j>0}|z_j|+r_t\|h\|+r_n\sqrt{N_J}\nonumber\\
 &\quad+\pi\left[
 m_y(KX+r)^2+
 \min\{\sqrt{M_z}\,X,K\sqrt{N_S}\,x_\infty\}X
 +2\frac{\sqrt{N_S}}{\sqrt6}Xr
 +\frac{\sqrt{M_S}}{\sqrt6}r^2
 \right].\label{round:c-single-support}
\end{align}
This is the necessary support bound from
Section~\ref{sec:exclusion}.
Its coefficient errors against ideal arrays satisfy
\[
 |\Delta a_y|<7\cdot10^{-12},\qquad
 |\Delta z_j|,|\Delta\widetilde z_j|<7\cdot10^{-12},\qquad
 |\Delta h_a|<4.3\cdot10^{-11}.
 \]
The first two use six input errors and a short dot.
The last uses six B errors and the same dot estimate.

The square roots in \eqref{round:c-single-support} are norms, so
they are stable also near zero.
For example $\sqrt{N_J}$ is the norm of the 36 components
$y_iJ_{ij}$; their perturbation norm is at most
$6\cdot10^{-12}$ before arithmetic, below $10^{-11}$ overall.
Similarly $\sqrt{N_S}$ is the norm of the complex components
$y_iS_i$, and $\sqrt{M_S}$ their maximum modulus; their errors are
below $10^{-11}$.
Finally $\sqrt{M_z}$ is the maximum norm of the two-component pairs
$(z_j,\widetilde z_j)$.
Each such norm changes by at most
$\sqrt2\,7\cdot10^{-12}<1.1\cdot10^{-11}$.
A maximum and a minimum of a fixed collection of approximate scalar
values have error at most the largest of their component errors;
no agreement of the selected index is needed.

For completeness, the linear h effect in
\eqref{round:c-single-support} is below
$.030001\sqrt{10}(4.3\cdot10^{-11})<4.1\cdot10^{-12}$;
the z effect is below $5a(7\cdot10^{-12})<3\cdot10^{-12}$;
the normal term contributes less than
$.019\cdot10^{-11}$.
The first Taylor-norm choice contributes less than
$\pi(.19)^2(1.1\cdot10^{-11})<1.3\cdot10^{-12}$,
and the alternative choice is smaller at the quoted scales.
The mixed and pure-r terms contribute less than
$2\pi(.19)(.03)10^{-11}/\sqrt6$ and
$\pi(.03)^2 10^{-11}/\sqrt6$.
The remaining positive-factor arithmetic is below $10^{-10}$,
so the entire possible support undererror is below $10^{-9}$.
Together with the constant-pairing error, this is well inside
the actual $2\cdot10^{-7}$ comparison slack.

\paragraph{Vector-ball radii.}
A rational bin midpoint is formed with error below
$5\cdot10^{-16}$ per variable.
All stored centers used by clustering are finite and of magnitude
at most two; this holds for new midpoint centers and is checked
after each update.
For a nonfar difference, wrapped subtraction and subtraction of
its mean have Euclidean error below $2\cdot10^{-14}$ against the
exact projection for the selected integer choices.
The six-component distance norm costs less than $10^{-13}$ more.
The last assignment therefore covers the entire bin by the radius
\[
 \sqrt{\mathtt{dst}}+\mathtt{layer.xnorm}+10^{-9}.
\]
For a newly inserted center, \texttt{dst} is zero and the midpoint
error alone is already covered.
The final assignment pass does not move its centers.
Intermediate averaging passes are only proposals, and the last
radii are recomputed from actual assignments.
Nonfinite or oversized centers, a failed radius check, or a cluster
cap cause failure, not an exclusion.

\subsection{Pair edges and the joint directional test}
\label{round:joint}

Let two accepted vector balls have radii $s_0,s_1<.1$.
The stored pair radius $r_e=s_0+s_1+10^{-10}$ is an upper bound
for their sum and satisfies $r_e<.201$.
Each product $Q_j$ of the two stored unit phases has complex error
below $1.3\cdot10^{-13}$, its six-term sum S below
$8\cdot10^{-13}$, and its squared modulus below
$1.1\cdot10^{-11}$; the last uses $|S|\le6$ in
\eqref{round:squared-perturbation}.
The edge constants and derivatives, for either O or U, have error
below $10^{-12}$ per real component.
For the U derivative, for instance, the product of S with a unit
coefficient is divided by six, leaving error at most
\[
 \frac{8\cdot10^{-13}+6(1.3\cdot10^{-13})}{6}
       +\text{short product rounding}<10^{-12}.
\]
The O constant is S divided by $2\pi$, and the U constant is
$(|S|^2/6-1)/(4\pi)$, which satisfy the same budget.

The edge-remainder bounds are evaluated with an added
$2\cdot10^{-8}$.
Their positive-operation error, including the input error in $|S|$,
is below $10^{-10}$; at $r_e<.201$ their uninflated values are
below $.255$.
The direct distance and support comparisons, including the O test
in direction S and the norm $\sqrt{\mathtt{jj}}$, have combined
possible error below $10^{-9}$.
Indeed the latter norm has coefficient perturbation at most
$\sqrt6\,10^{-12}$; its multiplier is at most $6r_e<1.206$,
and the constant/norm-S effects are at most the displayed
$1.1\cdot10^{-11}$ scale times bounded factors.
Their actual comparison slack is $2\cdot10^{-7}$.

For the joint test there are at most $6+6+2=14$ rows.
The proposed direction is checked finite and normalized as before,
so its stored entries have magnitude at most one.
The constant pairing has error below $1.5\cdot10^{-11}$.
Each accumulated shared-Hadamard coefficient has error below
$9\cdot10^{-11}$, since it receives at most twelve B input errors
of $7\cdot10^{-12}$ plus short-dot rounding.
Each vector-phase coefficient, including its edge contribution,
has error below $9\cdot10^{-12}$: six single-vector entries and at
most two edge entries each contribute at most $10^{-12}$.

Within each call to \texttt{single::support}, the vector-phase
coefficient array starts at zero.  The pure-vector remainder is
formed using its J sum and real companion sum before the joint
routine adds the edge derivatives.
Thus its maximum two-component coefficient norm has the same
$1.1\cdot10^{-11}$ error bound as in
Section~\ref{round:single}, with $X$ replaced by $s_t<.1$ and only
the first coefficient-norm alternative retained.
The weighting of the common B norm has input effect below
\[
 .030001\sqrt{10}\,9\cdot10^{-11}<8.6\cdot10^{-12}.
\]
Each vector-phase norm has perturbation at most
$\sqrt6\,9\cdot10^{-12}$, multiplied by $s_t<.1$.
The two single-vector residual supports have the norm-Lipschitz
errors just proved; the edge residual upper multiplies a direction
norm of at most $\sqrt2$ and has only ordinary bounded norm/product
rounding.  Adding these contributions and the positive scalar
arithmetic gives total possible support undererror below $10^{-8}$.
The actual comparison adds $3\cdot10^{-7}$.
Together with the constant-pairing error this certifies the joint
rejection.
No edge deletion depends on accuracy of the heuristic solves or
their weight updates.

\subsection{Finite paths and preservation of inconclusive cases}
\label{round:finite}

The preceding estimates apply to bounded finite certificate
evaluations.  We finish by explaining why exceptional heuristic
values cannot bypass those preconditions.
\begin{enumerate}
\item
In the seed cover, all direct lists come from finite integer
midpoints and \texttt{expair}.  The dual solver's entries must pass
the negated bounded comparison $|v_a|\le10^4$; NaN and infinity
fail it.  A failed proposal leaves the large safe support bound.
Gap/hash/reach cuts that are not necessary feasibility tests select
only possible reuse attempts.
\item
For A, a failed frame normalization makes setup fail, and the
entry point records an unresolved seed ball.  It does not certify
that seed ball empty.  The passed frames, fixed radius, and dyadic
unitaries give finite bounded block coefficients.  The potentially
ill-conditioned solves only propose directions; the final
\texttt{trial::certificate} checks every direction entry and its
maximum before use.
\item
The snapping iteration calls \texttt{gen}, which rejects
nonfinite or out-of-range phases.  Root certification requires
$\mathtt{val}<10^{-18}$, checks proposed eigenvalues finite,
checks HV/LV before using them, and checks T before its certificate
products.  Those bounded arrays make the products finite.
The optional tensor change of basis is range-checked; an unusable
one cannot lower the previously valid C bound.
The final conjunction of strict Root thresholds also rejects NaN.
\item
Candidate's inverse-angle output passes a finite range condition
before its exponential and a unit-direction comparison afterwards.
Its magnitude and chord guards bound every normalization.
Failure is an unsuccessful coverage proposal.
In \texttt{ell}, the data are range/symmetry checked, invalid
Cholesky pivots fail, and every final trial-solve coordinate is
checked before a residual certificate is evaluated.
The bounded residual scales exclude overflow.
The explicit sentinel values are outside the usable bootstrap
range or far above the target radius.
\item
B/C operate on genuine certified Records and also enforce the
stated numeric ranges in \texttt{validate}.
The read-time \texttt{finish} may form floating dephasings before
validation, but performs no conversion of those values to an
integer index; validation rejects invalid phases before they can
reach indexed geometric searches.
B's relation and grid formulas are then finite and bounded as
proved above.  A failed relation only removes a coverage
opportunity; a failed live-box solve keeps the parent unresolved.
\item
C's direct scalar and pair formulas have bounded finite inputs
and nonnegative square-sum arguments, so explicit finite checks at
every individual arithmetic operation are unnecessary.
Every potentially ill-conditioned directional solve is separately
checked and normalized.  Its failure keeps the bin or edge.
Clustering checks final center and radius bounds; caps and
incomplete graph searches report possibility/failure, as described
in Section~\ref{sec:exclusion}, rather than a negative certificate.
\end{enumerate}
In particular, \texttt{min}/\texttt{max} are not used to turn an
unchecked NaN from a heuristic solve into a valid small certificate.
Within successful bound evaluations their inputs are finite, and
a minimum of two established upper bounds remains an upper bound.
Within proposal generation, a bad proposal is checked before it can
justify an exclusion or cover.  These facts also explain why a
less successful heuristic or a different valid rounding outcome
may increase unresolved work without invalidating a successful
certificate.

\paragraph{Summary of the comparison margins.}
The following table collects the numerical separation used above;
\emph{error} means the entire adverse comparison error or the stated
local upper-bound error, not the error of an untrusted solver.
\begin{center}
\small
\begin{tabular}{p{.37\linewidth}p{.27\linewidth}p{.26\linewidth}}
\hline
Certificate path & Established adverse error & Explicit margin\\
\hline
seed gap / direct support & $10^{-14}$ / $10^{-11}$ & $10^{-8}$\\
shared-coordinate numerator & $<3.5\cdot10^{-11}$ & $10^{-7}$\\
shared-coordinate endpoints & $10^{-7}$ & $10^{-5}$\\
seed dual support & $10^{-4}$ & $.01$\\
A flat norm & $10^{-6}$ & $10^{-5}$\\
A directional comparison & $<4\cdot10^{-6}$ & $.0001$\\
Root identity Frobenius check & $<5\cdot10^{-8}$ & $2\cdot10^{-6}$\\
Root orthogonal / horizontal checks & $10^{-10}$ / $10^{-11}$ & $10^{-8}$\\
tensor transformed row bound & $10^{-6}$ & $.0002$\\
Candidate direct support & $2\cdot10^{-6}$ & $.0001$\\
Candidate external tangent error & $.00012$ radians & $.00015$ radians\\
first / second G entry error & $10^{-11}$ / $10^{-9}$ & $10^{-7}I$\\
spectral row / fourth-norm bound & $3\cdot10^{-13}$ / $10^{-10}$ &
 $10^{-10}$ / $2\cdot10^{-7}$\\
bootstrap update & $10^{-12}$ & $10^{-8}$\\
B projected normal / local update & $5\cdot10^{-11}$ / $10^{-13}$ & $10^{-8}$\\
B cross-Gram spectral perturbation & $2\cdot10^{-9}$ & $10^{-7}$ before root\\
C direct / directional comparison & $10^{-9}$ & $2\cdot10^{-7}$\\
C joint support & $10^{-8}$ & $3\cdot10^{-7}$\\
\hline
\end{tabular}
\end{center}
The resolvent signed-addition argument is given separately in
Section~\ref{round:ell}, since its absolute scale can be large and
its proof uses both an absolute dot surplus and a relative quotient
surplus.  All margins in the table are used in the source at the
places proved above; none is inferred merely from an observed
distance from a threshold.
